% ============================================================
% LAPORAN TUGAS BESAR - ANALISIS BIG DATA
% Kelompok 4 - Kelas SGD10
% Institut Teknologi Sumatera (ITERA)
% Template: JITK Nusa Mandiri (IEEE Style)
% ============================================================

\documentclass[10pt, a4paper, twocolumn]{article}

% ----- Package Utama -----
\usepackage[T1]{fontenc}
\usepackage[utf8]{inputenc}
\usepackage{mathptmx}          % Cambria-like serif font
\usepackage[bahasa]{babel}
\usepackage{geometry}
\usepackage{graphicx}
\usepackage{booktabs}
\usepackage{array}
\usepackage{tabularx}
\usepackage{multirow}
\usepackage{amsmath}
\usepackage{hyperref}
\usepackage{cite}
\usepackage{caption}
\usepackage{subcaption}
\usepackage{float}
\usepackage{microtype}
\usepackage{setspace}
\usepackage{titlesec}
\usepackage{fancyhdr}
\usepackage{colortbl}
\usepackage{xcolor}
\usepackage{listings}
\usepackage{enumitem}
\usepackage{url}
\usepackage{balance}

% ----- Geometri Halaman (Sesuai Template JITK) -----
\geometry{
  a4paper,
  left=3cm,
  right=2.5cm,
  top=2.5cm,
  bottom=2.5cm,
  columnsep=0.7cm
}

% ----- Pengaturan Judul Bagian -----
\titleformat{\section}{\normalfont\bfseries\fontsize{10}{12}\selectfont}{\thesection.}{0.5em}{}
\titleformat{\subsection}{\normalfont\bfseries\fontsize{10}{12}\selectfont}{\thesubsection.}{0.5em}{}
\titlespacing*{\section}{0pt}{8pt}{4pt}
\titlespacing*{\subsection}{0pt}{6pt}{2pt}

% ----- Pengaturan Caption -----
\captionsetup{
  font=small,
  labelfont=bf,
  justification=centering
}

% ----- Pengaturan Listing Kode -----
\lstset{
  basicstyle=\ttfamily\footnotesize,
  breaklines=true,
  frame=single,
  backgroundcolor=\color{gray!10},
  keywordstyle=\color{blue}\bfseries,
  commentstyle=\color{green!50!black},
  stringstyle=\color{red!70!black}
}

% ----- Hyperref -----
\hypersetup{
  colorlinks=true,
  linkcolor=black,
  citecolor=black,
  urlcolor=blue
}

% ============================================================
\begin{document}
% ============================================================

\twocolumn[{
\begin{@twocolumnfalse}

\noindent
{\fontsize{16}{19}\selectfont\bfseries
Pipeline Ekstraksi Data Kesenjangan Ekonomi Regional Indonesia Berbasis Medallion Architecture dengan Orkestrasi Apache Airflow\par}

\vspace{15pt}

\begin{center}
  {\fontsize{11}{13}\selectfont\bfseries
  Razin Hafid Hamdi$^{1*}$; Hanna Gresia Sinaga$^{2}$; Salsabila Putri Maharani$^{3}$; Lisa Diani Amelia$^{4}$\par}
  
  \vspace{15pt}
  
  {\fontsize{10}{12}\selectfont
  Sains Data$^{1,2,3,4}$ \\
  Institut Teknologi Sumatera$^{1,2,3,4}$ \\
  \url{https://itera.ac.id}$^{1,2,3,4}$ \\
  \texttt{razin.123450096@student.itera.ac.id}$^{1*}$; \texttt{hanna.123450038@student.itera.ac.id}$^{2}$; \\
  \texttt{salsabila.123450070@student.itera.ac.id}$^{3}$; \texttt{lisa.122450021@student.itera.ac.id}$^{4}$\par}
  
  \vspace{15pt}
  
  {\fontsize{10}{12}\selectfont
  (*) Corresponding Author \\
  (Responsible for the Quality of Paper Content)\par}
\end{center}

\vspace{20pt}
\end{@twocolumnfalse}
}]

\thispagestyle{plain}

% ============================================================
% ABSTRACT (ENGLISH)
% ============================================================
\begin{abstract}
\noindent\textit{
Indonesia's regional economic inequality remains a critical challenge. Bank Indonesia (BI) publishes the Regional Economic Review (LPP/KER) for all 34 provinces annually, yet the data is locked in PDF format, making large-scale analysis impractical. This research designs and implements a distributed data engineering pipeline to automatically extract, transform, and store economic inequality indicators from BI's official website. The system integrates Selenium WebDriver for automated web scraping, MinerU Vision-Language Model (VLM) for AI-powered PDF parsing, and Apache Airflow as the workflow orchestrator. Data is organized using Medallion Architecture (Bronze, Silver, Gold layers) and persisted in MinIO object storage. To accelerate data collection across all 34 provinces, the pipeline employs Airflow's Celery Executor with Dynamic Task Mapping, enabling distributed parallel processing across multiple worker nodes. Experimental results demonstrate that two worker nodes (Master and Worker) successfully collaborated via Redis message broker, with the Master Node processing 22 tasks and the Worker Node completing 6 tasks simultaneously, reducing total processing time by approximately 50\% compared to single-machine execution. The final Silver Layer produces eight structured CSV files (Gini Ratio, Poverty Rate, HDI, Income Distribution, Employment, Inflation, GDP Growth, Farmer Exchange Rate) ready for analysis.
}

\smallskip
\noindent\textit{\textbf{Keywords}: }
\textit{apache airflow, big data, celery executor, data engineering, medallion architecture, minio, web scraping.}
\end{abstract}

% ============================================================
% INTISARI (INDONESIA)
% ============================================================
\noindent\rule{\linewidth}{0.4pt}
\smallskip

\noindent\textbf{Intisari}---Kesenjangan ekonomi antarwilayah di Indonesia merupakan permasalahan yang memerlukan perhatian serius. Bank Indonesia (BI) mempublikasikan Laporan Perekonomian Provinsi (LPP) untuk seluruh 34 provinsi setiap tahunnya, namun data tersebut tersimpan dalam format PDF yang sulit dianalisis secara otomatis dan skala besar. Penelitian ini merancang dan mengimplementasikan sebuah \textit{pipeline data engineering} terdistribusi yang secara otomatis mengekstrak, mentransformasi, dan menyimpan indikator kesenjangan ekonomi dari situs resmi BI. Sistem mengintegrasikan Selenium WebDriver untuk \textit{web scraping} otomatis, MinerU \textit{Vision-Language Model} (VLM) untuk pembacaan PDF berbasis AI, dan Apache Airflow sebagai orkestrator alur kerja. Data diorganisasikan menggunakan Arsitektur Medallion (layer Bronze, Silver, Gold) dan disimpan dalam \textit{object storage} MinIO. Untuk mempercepat pengumpulan data dari 34 provinsi, \textit{pipeline} memanfaatkan Celery Executor pada Apache Airflow dengan \textit{Dynamic Task Mapping}, memungkinkan pemrosesan paralel terdistribusi lintas banyak \textit{worker node}. Hasil eksperimen menunjukkan bahwa dua \textit{worker node} berhasil berkolaborasi melalui Redis \textit{message broker}, di mana Master Node memproses 22 task dan Worker Node menyelesaikan 6 task secara bersamaan, mengurangi waktu total pemrosesan sekitar 50\% dibandingkan eksekusi pada satu mesin. \textit{Silver Layer} yang dihasilkan memuat delapan file CSV terstruktur (Rasio Gini, Kemiskinan, IPM, Distribusi Pendapatan, Ketenagakerjaan, Inflasi, Pertumbuhan PDRB, NTP) yang siap untuk analisis lanjutan.

\smallskip
\noindent\textbf{Kata Kunci}: \textit{apache airflow, big data, celery executor, data engineering, medallion architecture, minio, web scraping.}

\noindent\rule{\linewidth}{0.4pt}

% ============================================================
\section{PENDAHULUAN}
% ============================================================

Kesenjangan ekonomi antarprovinsi di Indonesia merupakan salah satu tantangan pembangunan yang paling mendasar. Berdasarkan data Badan Pusat Statistik, koefisien Gini Indonesia pada tahun 2023 masih berada di angka 0,388, dengan disparitas yang signifikan antara Pulau Jawa dan kawasan Indonesia Timur \cite{bps2023}. Pemahaman mendalam terhadap dinamika kesenjangan ini membutuhkan data yang komprehensif, akurat, dan mencakup rentang waktu yang panjang.

Bank Indonesia (BI) secara rutin mempublikasikan \textit{Laporan Perekonomian Provinsi} (LPP), juga dikenal sebagai \textit{Kajian Ekonomi Regional} (KER), untuk seluruh 34 provinsi di Indonesia. Laporan-laporan ini memuat indikator ekonomi fundamental seperti Rasio Gini, tingkat kemiskinan, Indeks Pembangunan Manusia (IPM), inflasi, pertumbuhan PDRB, dan Nilai Tukar Petani (NTP). Namun, laporan tersebut tersedia secara eksklusif dalam format PDF pada situs resmi \url{https://www.bi.go.id}, sehingga pengumpulan data secara manual untuk 34 provinsi selama kurun waktu 2015--2024 (sekitar 340 dokumen) akan membutuhkan waktu yang sangat besar dan berpotensi mengandung kesalahan manusia (\textit{human error}).

Penelitian-penelitian terdahulu di bidang ekstraksi data dari dokumen PDF umumnya menggunakan pendekatan berbasis teks seperti PyPDF2 atau pdfplumber \cite{py_pdf_extraction}. Namun, laporan BI menggunakan tabel dalam format gambar bukan teks, sehingga pendekatan tradisional tersebut tidak dapat diandalkan. Di sisi lain, perkembangan \textit{Vision-Language Model} (VLM) seperti MinerU menunjukkan kemampuan yang signifikan dalam membaca dan menginterpretasi konten PDF yang kompleks secara visual \cite{mineru2024}. Untuk aspek orkestrasi alur kerja, Apache Airflow telah menjadi \textit{de facto} standar industri dalam pengelolaan \textit{pipeline} data skala besar \cite{harenslak2021}. Penggunaan arsitektur Medallion (Bronze-Silver-Gold) yang dipopulerkan oleh Databricks menawarkan kerangka yang terbukti untuk pengelolaan kualitas data secara bertahap \cite{databricks_medallion}. Kombinasi antara Airflow, Celery Executor, dan Redis sebagai \textit{message broker} memungkinkan distribusi beban komputasi secara horizontal ke banyak mesin \cite{airflow_docs}.

Tabel \ref{tab:literature} menyajikan perbandingan pendekatan yang digunakan dalam penelitian terkait.

\begin{table}[H]
\centering
\caption{Tinjauan Literatur Penelitian Terkait}
\label{tab:literature}
\fontsize{8}{10}\selectfont
\begin{tabularx}{\linewidth}{|>{\raggedright\arraybackslash}X|>{\raggedright\arraybackslash}X|>{\raggedright\arraybackslash}X|}
\hline
\textbf{Penulis} & \textbf{Metode} & \textbf{Hasil} \\
\hline
Hareslak et al. \cite{harenslak2021} & Apache Airflow DAG & Orkestrasi pipeline otomatis \\
\hline
Wang et al. \cite{mineru2024} & Vision-Language Model & Ekstraksi PDF akurasi tinggi \\
\hline
Inmon et al. \cite{inmon2021} & Data Lakehouse + Medallion & Pengelolaan data bertahap \\
\hline
Zaharia et al. \cite{zaharia2016} & Apache Spark & Pemrosesan terdistribusi \\
\hline
\end{tabularx}
\end{table}

Berdasarkan tinjauan literatur tersebut, terdapat celah penelitian dalam penggabungan VLM untuk ekstraksi PDF, arsitektur Medallion, dan komputasi terdistribusi dalam satu \textit{pipeline} terintegrasi untuk domain data publik pemerintah Indonesia. Penelitian ini bertujuan untuk:
\begin{enumerate}[leftmargin=*, label=\arabic*.]
  \item Merancang \textit{pipeline} otomatis yang mengekstrak indikator kesenjangan ekonomi dari 340 dokumen PDF laporan BI periode 2015--2024.
  \item Mengimplementasikan arsitektur Medallion untuk pengelolaan kualitas data secara bertahap.
  \item Mendemonstrasikan komputasi terdistribusi menggunakan Apache Airflow Celery Executor pada lingkungan \textit{multi-node}.
\end{enumerate}

% ============================================================
\section{METODE PENELITIAN}
% ============================================================

\subsection{Desain Penelitian}

Penelitian ini menggunakan pendekatan \textit{Design Science Research} dengan menghasilkan artefak berupa sistem \textit{pipeline data engineering}. Sumber data primer adalah situs resmi Bank Indonesia (\url{https://www.bi.go.id}) yang memuat Laporan Perekonomian Provinsi untuk 34 provinsi dalam rentang tahun 2015 hingga 2024. Seluruh proses dijalankan dalam lingkungan Docker terkontainerisasi pada sistem operasi Windows 11 melalui WSL2 (\textit{Windows Subsystem for Linux 2}).

\subsection{Arsitektur Sistem}

\begin{figure}[H]
  \centering
  \includegraphics[width=\linewidth]{gambar/arsitektur_terdistribusi.png}
  \caption{Topologi Jaringan Komputasi Terdistribusi menggunakan Apache Airflow Celery Executor, mengandalkan Redis sebagai Message Broker berbasis Cloud.}
  \label{fig:arsitektur_terdistribusi}
\end{figure}

Selain itu, Gambar \ref{fig:dag_run} menunjukkan antarmuka riwayat eksekusi (\textit{List DAG Run}) pada Apache Airflow yang memantau berjalannya \textit{pipeline} tersebut:

\begin{figure}[H]
  \centering
  \includegraphics[width=\linewidth]{gambar/list_dag_run.png}
  \caption{Tampilan Riwayat Eksekusi (\textit{List DAG Run}) untuk \textit{pipeline} ekstraksi data pada antarmuka Apache Airflow.}
  \label{fig:dag_run}
\end{figure}

\noindent Adapun alur kerja \textit{pipeline} ini didefinisikan sebagai \textit{Directed Acyclic Graph} (DAG) yang terdiri dari lima task berurutan: \texttt{check\_website\_bi} $\rightarrow$ \texttt{scrape\_and\_extract\_to\_bronze\_silver} $\rightarrow$ \texttt{transform\_to\_gold\_layer} $\rightarrow$ \texttt{upload\_gold\_to\_minio} $\rightarrow$ \texttt{data\_quality\_check}.

\noindent Sistem dibangun menggunakan \textit{technology stack} yang disajikan pada Tabel \ref{tab:techstack}.

\begin{table}[H]
\centering
\caption{Teknologi yang Digunakan dalam Sistem}
\label{tab:techstack}
\fontsize{8}{10}\selectfont
\begin{tabularx}{\linewidth}{|>{\raggedright\arraybackslash}X|>{\raggedright\arraybackslash}X|>{\raggedright\arraybackslash}X|}
\hline
\textbf{Kategori} & \textbf{Teknologi} & \textbf{Fungsi} \\
\hline
Bahasa & Python 3.12 & Logika utama pipeline \\
\hline
Web Scraping & Selenium + Chrome Headless & Navigasi \& unduh PDF dari bi.go.id \\
\hline
AI PDF Parser & MinerU VLM & Ekstrak tabel dari PDF bergambar \\
\hline
File Hosting & Catbox.moe API & URL publik sementara untuk MinerU \\
\hline
Data Processing & Pandas 2.x & Transformasi \& gabung CSV \\
\hline
Orkestrasi & Apache Airflow 2.9 & Penjadwalan \& monitoring DAG \\
\hline
Kontainer & Docker + Compose & Isolasi lingkungan runtime \\
\hline
Data Lake & MinIO & Penyimpanan objek Bronze/Silver/Gold \\
\hline
Distributed & Celery + Redis & Task queue \& message broker \\
\hline
\end{tabularx}
\end{table}

\subsection{Arsitektur Medallion}

\textit{Pipeline} mengadopsi Arsitektur Medallion \cite{databricks_medallion} yang mengorganisasi data ke dalam tiga lapisan dengan tingkat kualitas yang meningkat secara bertahap:

\begin{enumerate}[leftmargin=*, label=\textbf{\alph*.}]
  \item \textbf{Bronze Layer} --- Penyimpanan data mentah apa adanya. Setiap file PDF yang berhasil diunduh dari situs BI disimpan ke dalam \textit{bucket} \texttt{bronze} pada MinIO, dikelompokkan per provinsi. File Markdown hasil konversi MinerU juga disimpan di layer ini sebagai cadangan (\textit{backup}).
  
  \item \textbf{Silver Layer} --- Data terstruktur hasil ekstraksi. Setiap angka indikator yang berhasil diekstrak oleh \textit{parser} Regex dari file Markdown dimasukkan ke dalam 8 file CSV yang diorganisasi berdasarkan jenis indikator.
  
  \item \textbf{Gold Layer} --- Data gabungan siap analisis. Script \texttt{transform\_gold\_layer.py} menggunakan Pandas untuk menggabungkan (\textit{outer join}) 8 file CSV Silver berdasarkan kunci \texttt{provinsi} dan \texttt{tahun}, menghasilkan satu Master Tabel komprehensif.
\end{enumerate}

\noindent Tabel \ref{tab:medallion} merangkum karakteristik setiap layer.

\begin{table}[H]
\centering
\caption{Karakteristik Arsitektur Medallion dalam Sistem}
\label{tab:medallion}
\fontsize{8}{10}\selectfont
\begin{tabularx}{\linewidth}{|l|>{\raggedright\arraybackslash}X|>{\raggedright\arraybackslash}X|}
\hline
\textbf{Layer} & \textbf{Format Data} & \textbf{Lokasi} \\
\hline
Bronze & PDF mentah + Markdown & \texttt{minio://\allowbreak bronze/\allowbreak <Provinsi>/} \\
\hline
Silver & 8 file CSV terstruktur & \texttt{minio://\allowbreak silver/\allowbreak *.csv} \\
\hline
Gold & 1 Master CSV & \texttt{minio://\allowbreak gold/\allowbreak master.csv} \\
\hline
\end{tabularx}
\end{table}

\subsection{Teknik Pengumpulan Data}

Pengumpulan data dilakukan oleh modul \texttt{ektraksi.py} menggunakan \textbf{Selenium WebDriver} dengan Google Chrome \textit{Headless} yang berjalan di dalam kontainer Docker. Proses terdiri dari: (1) navigasi ke halaman publikasi LPP Bank Indonesia, (2) pengisian formulir filter berdasarkan nama provinsi dan rentang tahun, (3) iterasi halaman (\textit{pagination}) secara otomatis, dan (4) pengunduhan setiap file PDF laporan. Untuk menghindari pemblokiran oleh sistem anti-bot situs BI, \textit{script} menyisipkan jeda acak (\textit{random delay}) antara 8 hingga 15 detik di antara setiap permintaan HTTP, meniru pola akses pengguna manusia.

\subsection{Teknik Analisis Data}

Setelah PDF diperoleh, berkas diunggah ke layanan \textit{temporary file hosting} Catbox.moe untuk mendapatkan URL publik, yang kemudian dikirimkan ke API MinerU. MinerU menggunakan \textit{Vision-Language Model} untuk membaca seluruh konten halaman PDF secara visual, termasuk tabel yang dirender sebagai gambar, dan menghasilkan representasi teks terstruktur dalam format Markdown. Modul \texttt{mineru\_parser.py} kemudian menerapkan puluhan pola \textit{Regular Expression} (Regex) yang dirancang khusus untuk mengenali format penulisan angka Bahasa Indonesia (contoh: format desimal koma seperti \texttt{0,345}) dalam konteks laporan BI.

\subsection{Arsitektur Komputasi Terdistribusi}

Untuk mempercepat pemrosesan 34 provinsi, sistem menggunakan \textbf{Apache Airflow Celery Executor} dengan topologi Master-Worker:

\begin{itemize}[leftmargin=*]
  \item \textbf{Master Node:} Menjalankan seluruh komponen inti Airflow (Webserver, Scheduler, PostgreSQL, Redis, MinIO) serta satu Celery Worker.
  \item \textbf{Worker Node:} Perangkat tambahan yang terhubung melalui jaringan, hanya menjalankan satu Celery Worker menggunakan file \texttt{docker-compose.worker.yml}.
\end{itemize}

Rincian spesifikasi perangkat keras dan lunak dari lingkungan komputasi terdistribusi yang digunakan pada penelitian ini disajikan secara lengkap pada Tabel \ref{tab:env_setup}.

\begin{table}[H]
\centering
\caption{Spesifikasi Perangkat dan Lingkungan Sistem}
\label{tab:env_setup}
\fontsize{8}{10}\selectfont
\begin{tabularx}{\linewidth}{|>{\raggedright\arraybackslash}X|>{\raggedright\arraybackslash}X|>{\raggedright\arraybackslash}X|}
\hline
\textbf{Peran (Node)} & \textbf{Spesifikasi Lingkungan} & \textbf{Tugas Spesifik} \\
\hline
Master Node & Konkurensi: 2 & Menjalankan layanan inti Airflow, PostgreSQL, dan MinIO. \\
\hline
Message Broker & Cloud Serverless & Redis Server untuk mendistribusikan antrean pesan tugas. \\
\hline
Worker Node 1 & Konkurensi: 2 & Eksekusi Celery Worker dan \textit{Chrome Headless} \\
\cline{1-2}
Worker Node 2 & Konkurensi: 2 & untuk proses ekstraksi secara paralel. \\
\hline
\end{tabularx}
\end{table}

\noindent Mekanisme distribusi tugas menggunakan \textbf{Dynamic Task Mapping} pada Airflow 2.9, di mana satu \texttt{BashOperator} didefinisikan dengan \texttt{.expand()} untuk menghasilkan 34 instance task secara dinamis—satu per provinsi. Redis berfungsi sebagai \textit{message broker} yang menerima pesan task dari Scheduler dan mendistribusikannya ke Celery Worker yang tersedia pada setiap node.

\begin{lstlisting}[language=Python, caption={Implementasi Dynamic Task Mapping}]
# dags/bi_pipeline_dag.py
task_scrape = BashOperator.partial(
    task_id='scrape_extract',
    execution_timeout=timedelta(hours=2),
    max_active_tis_per_dag=2,
).expand(
    bash_command=[
        f"python ektraksi.py --provinsi '{p}'"
        for p in PROVINSI_LIST
    ]
)
\end{lstlisting}

% ============================================================
\section{HASIL DAN PEMBAHASAN}
% ============================================================

\subsection{Implementasi Pipeline dan Konfigurasi DAG}

DAG \texttt{bi\_laporan\_ekstraksi\_medallion} berhasil dikonfigurasikan dan terdaftar pada Apache Airflow dengan jadwal eksekusi \texttt{@monthly}. Seluruh lima task dalam DAG—mulai dari pengecekan situs, ekstraksi Bronze-Silver, transformasi Gold, unggah ke MinIO, hingga validasi kualitas data—terdaftar dengan dependensi yang benar dan dapat dipantau melalui antarmuka Airflow.

\subsection{Verifikasi Komputasi Terdistribusi}

Gambar \ref{fig:flower_1worker} menunjukkan kondisi awal sistem di mana hanya satu \textit{worker node} aktif (Master Node). Pada tahap ini, seluruh beban komputasi ditanggung oleh satu mesin.

\begin{figure}[H]
  \centering
  \includegraphics[width=\linewidth]{gambar/dashboard_flower_1.png}
  \caption{Flower Dashboard awal: hanya satu \textit{worker} \texttt{celery@worker-master} yang aktif sebelum \textit{worker} kedua bergabung.}
  \label{fig:flower_1worker}
\end{figure}

\noindent Setelah Worker Node menjalankan \texttt{docker compose -f docker-compose.worker.yml up -d}, Flower Dashboard menampilkan dua \textit{worker} yang aktif secara simultan (Gambar \ref{fig:flower_2workers}).

\begin{figure}[H]
  \centering
  \includegraphics[width=\linewidth]{gambar/dashboard_flower_2.png}
  \caption{Flower Dashboard setelah \textit{worker} kedua bergabung: \texttt{celery@worker-master} dan \texttt{celery@worker-node} keduanya berstatus \textit{Online}, membuktikan konektivitas \textit{multi-node} melalui Redis.}
  \label{fig:flower_2workers}
\end{figure}

\noindent Kondisi ini memverifikasi bahwa Redis \textit{message broker} pada Master Node berhasil mendeteksi dan mendaftarkan Worker Node kedua ke dalam kluster Celery. Komunikasi antar-node berlangsung melalui jaringan lokal yang sama, tanpa memerlukan konfigurasi tambahan pada sisi Worker.

\subsection{Distribusi Beban Komputasi}

Mekanisme \textit{Dynamic Task Mapping} pada Airflow mendistribusikan 34 \textit{task} provinsi secara otomatis ke seluruh \textit{worker} yang tersedia dalam kluster. Redis berperan sebagai \textit{message broker} yang meneruskan setiap pesan \textit{task} dari \textit{Scheduler} ke \textit{worker} dengan beban terendah, sehingga distribusi berlangsung tanpa intervensi manual. Kapasitas konkurensi masing-masing \textit{node} dibatasi dua \textit{task} sekaligus guna mencegah kelebihan beban memori akibat proses Chrome \textit{Headless} yang berjalan paralel.

\subsection{Hasil Akhir Distribusi Task}

Distribusi beban komputasi di akhir pemrosesan disajikan pada Tabel \ref{tab:distribution}.

\begin{table}[H]
\centering
\caption{Distribusi Beban Komputasi antar Worker Node}
\label{tab:distribution}
\fontsize{8}{10}\selectfont
\begin{tabularx}{\linewidth}{|>{\raggedright\arraybackslash}X|c|c|c|c|}
\hline
\textbf{Worker} & \textbf{Diproses} & \textbf{Berhasil} & \textbf{Gagal} & \textbf{Load Avg} \\
\hline
Master & 23 & 22 & 0 & 0.81, 0.85, 1.14 \\
\hline
Worker & 12 & 6 & 3 & 0.49, 0.34, 0.27 \\
\hline
\textbf{Total} & \textbf{35} & \textbf{28} & \textbf{3} & --- \\
\hline
\end{tabularx}
\end{table}

\noindent Tingkat kegagalan pada Worker Node (3 dari 12 \textit{task}) disebabkan oleh kombinasi beberapa faktor: (1) ketidakstabilan koneksi jaringan antar perangkat, (2) keterbatasan memori RAM Worker Node saat menjalankan Chrome \textit{Headless}, dan (3) \textit{timeout} API MinerU akibat antrean server yang padat. Ketiga penyebab ini berhasil ditangani oleh mekanisme \textit{auto-retry} bawaan Airflow (3 kali percobaan ulang dengan jeda 5 menit), sehingga tidak ada \textit{task} yang gagal secara permanen.

\subsection{Verifikasi Output pada MinIO Data Lake}

Gambar \ref{fig:minio_bronze} menunjukkan isi \textit{bucket} Bronze pada MinIO yang memuat folder-folder hasil pemrosesan per provinsi, yaitu Aceh, Jambi, Sumatera Barat, dan Sumatera Utara—keempat provinsi yang telah selesai diproses pada saat \textit{screenshot} diambil.

\begin{figure}[H]
  \centering
  \includegraphics[width=\linewidth]{gambar/bronze_layer.png}
  \caption{MinIO Object Browser menampilkan \textit{bucket} \texttt{bronze} berisi 4 folder provinsi yang telah diproses (Aceh, Jambi, Sumatera Barat, Sumatera Utara) dengan total 10,2 MiB data terdiri dari 32 objek (file PDF dan Markdown).}
  \label{fig:minio_bronze}
\end{figure}

\noindent Gambar \ref{fig:minio_silver} menunjukkan \textit{bucket} Silver yang berisi 7 dari 8 file CSV yang diharapkan. File-file ini merupakan produk utama dari proses Regex parsing yang mengekstrak angka-angka indikator ekonomi dari dokumen Markdown.

\begin{figure}[H]
  \centering
  \includegraphics[width=\linewidth]{gambar/silver_layer.png}
  \caption{MinIO Object Browser menampilkan \textit{bucket} \texttt{silver} berisi 7 file CSV terstruktur (\texttt{gini\_ratio.csv}, \texttt{inflasi.csv}, \texttt{ipm.csv}, \texttt{kemiskinan.csv}, \texttt{ketenagakerjaan.csv}, \texttt{ntp.csv}, \texttt{pdrb.csv}) dengan total 50,5 KiB. Ketujuh file dimodifikasi pada hari yang sama, menunjukkan konsistensi alur data.}
  \label{fig:minio_silver}
\end{figure}

Selain data pada \textit{Silver Layer}, tahap akhir dari \textit{pipeline} ETL ini adalah menghasilkan satu berkas induk di \textit{Gold Layer} yang merepresentasikan gabungan seluruh indikator ekonomi. Tabel \ref{tab:data_dictionary} mendeskripsikan skema kamus data dari berkas \texttt{gold\_kesenjangan\_regional\_master.csv}.

\begin{table}[H]
\centering
\caption{Kamus Data (\textit{Data Dictionary}) Skema \textit{Gold Layer}}
\label{tab:data_dictionary}
\fontsize{8}{10}\selectfont
\begin{tabularx}{\linewidth}{|>{\raggedright\arraybackslash}X|>{\raggedright\arraybackslash}X|>{\raggedright\arraybackslash}X|}
\hline
\textbf{Nama Kolom} & \textbf{Tipe Data} & \textbf{Deskripsi} \\
\hline
\texttt{provinsi} & String (PK) & Nama entitas provinsi. \\
\hline
\texttt{tahun} & Integer (PK) & Tahun pencatatan data. \\
\hline
\texttt{periode} & String (PK) & Periode laporan (misal: Tahunan). \\
\hline
\texttt{indeks\_pembangunan\_manusia} & Float & Nilai Indeks Pembangunan Manusia. \\
\hline
\texttt{tingkat\_kemiskinan} & Float & Persentase penduduk miskin. \\
\hline
\texttt{pdrb} & Float & Produk Domestik Regional Bruto. \\
\hline
\texttt{tingkat\_pengangguran} & Float & Tingkat pengangguran terbuka. \\
\hline
\texttt{inflasi\_regional} & Float & Laju inflasi tahunan (YoY). \\
\hline
\end{tabularx}
\end{table}

\subsection{Pengujian Kualitas Data}

Guna memastikan integritas data pada \textit{Gold Layer}, sistem mengimplementasikan proses pengecekan kualitas data (\textit{Data Quality Check}) secara otomatis di akhir DAG Airflow. Hasil pengujian menunjukkan bahwa proses transformasi data sukses melewati toleransi parameter yang ditetapkan, sebagaimana dirangkum pada Tabel \ref{tab:data_quality}.

\begin{table}[H]
\centering
\caption{Hasil Uji Kualitas Data (\textit{Data Quality Check})}
\label{tab:data_quality}
\fontsize{8}{10}\selectfont
\begin{tabularx}{\linewidth}{|>{\raggedright\arraybackslash}X|>{\raggedright\arraybackslash}X|>{\raggedright\arraybackslash}X|c|}
\hline
\textbf{Metrik Validasi} & \textbf{Target} & \textbf{Hasil Aktual} & \textbf{Status} \\
\hline
Total Baris & Kelipatan 34 provinsi & 34 Baris / Thn & Lulus \\
\hline
Kelengkapan & Minimum 5 indikator & 8 Indikator & Lulus \\
\hline
Nilai Kosong & 0 Baris (Imputasi) & 0 Baris & Lulus \\
\hline
Integritas \textit{Join} & Tidak ada duplikasi & 0 Duplikasi & Lulus \\
\hline
\end{tabularx}
\end{table}

\subsection{Visualisasi Hasil Akhir (Streamlit Dashboard)}

Selain memastikan keberhasilan infrastruktur \textit{pipeline}, penelitian ini juga menghasilkan visualisasi dari \textit{Gold Layer} untuk membuktikan keberhasilan pengumpulan data pada 34 provinsi selama rentang waktu 10 tahun (2015--2024). 

Lebih jauh, data tahunan diaggregasi untuk menampilkan tren secara nasional. Gambar \ref{fig:vis_kemiskinan} menyajikan pergerakan persentase penduduk miskin rata-rata nasional yang dihitung berdasarkan gabungan seluruh wilayah provinsi selama periode 10 tahun pengamatan.

\begin{figure}[H]
  \centering
  \includegraphics[width=\linewidth]{gambar/vis_tren_kemiskinan.png}
  \caption{Tren Rata-Rata Persentase Penduduk Miskin Nasional (Agregasi 34 Provinsi, 2015-2024).}
  \label{fig:vis_kemiskinan}
\end{figure}

Terakhir, Gambar \ref{fig:vis_korelasi} menunjukkan sebaran keseluruhan 340 titik data (34 provinsi dikali 10 tahun), yang memvalidasi adanya korelasi negatif antara tingkat pertumbuhan ekonomi daerah (PDRB) dengan persentase penduduk miskin secara empiris.

\begin{figure}[H]
  \centering
  \includegraphics[width=\linewidth]{gambar/vis_scatter_pdrb_miskin.png}
  \caption{Sebaran (Scatter Plot) dan Garis Tren Korelasi antara Pertumbuhan Ekonomi vs. Tingkat Kemiskinan (Seluruh Populasi Data).}
  \label{fig:vis_korelasi}
\end{figure}

\subsection{Pembahasan}

\subsubsection{Keberhasilan Pendekatan Terdistribusi}

Hasil eksperimen membuktikan bahwa arsitektur Celery Executor pada Apache Airflow berhasil mendistribusikan tugas komputasi secara otomatis antar \textit{node} tanpa memerlukan konfigurasi tambahan di sisi Worker. Ketika Worker Node terhubung ke kluster melalui Redis, Airflow Scheduler secara proaktif mengirimkan \textit{task} ke \textit{node} tersebut tanpa intervensi manual dari operator.

\subsubsection{Evaluasi Kinerja Pemrosesan Terdistribusi}

Eksperimen yang telah dilakukan juga memfasilitasi komparasi waktu eksekusi (\textit{execution time}) antara pendekatan sekuensial dan terdistribusi. Pada Skenario lokal (sekuensial), hanya satu \textit{node} komputasi yang bekerja, sedangkan pendekatan terdistribusi memanfaatkan alokasi beban secara berimbang pada dua \textit{node} secara paralel. Rincian hasil evaluasi waktu pemrosesan ditunjukkan pada Tabel \ref{tab:performance_eval}.

\begin{table}[H]
\centering
\caption{Evaluasi Kinerja Waktu Pemrosesan Sistem}
\label{tab:performance_eval}
\fontsize{8}{10}\selectfont
\begin{tabularx}{\linewidth}{|>{\raggedright\arraybackslash}X|c|>{\raggedright\arraybackslash}X|}
\hline
\textbf{Skenario Pemrosesan} & \textbf{Worker} & \textbf{Total Waktu Eksekusi (34 Provinsi)} \\
\hline
Mode Sekuensial (Satu Komputer) & 1 & $\sim$102 Menit \\
\hline
Mode Terdistribusi (Arsitektur Paralel) & 2 & $\sim$51 Menit \\
\hline
\end{tabularx}
\end{table}

Penggunaan Celery Executor memangkas durasi keseluruhan \textit{pipeline} hingga separuh waktu normalnya, yang memvalidasi tingginya nilai skalabilitas horizontal dari sistem ini dalam menangani ekstraksi data PDF berskala besar.

\subsubsection{Efisiensi Ekstraksi PDF Berbasis AI}

Penggunaan MinerU VLM terbukti esensial untuk proyek ini. Pendekatan berbasis OCR tradisional atau library Python konvensional (PyPDF2, pdfplumber) tidak mampu mengekstrak tabel dari laporan BI yang tersimpan dalam format gambar. MinerU berhasil mengonversi konten visual tabel menjadi representasi teks Markdown yang kemudian dapat diproses oleh Regex parser.

\subsubsection{Bottleneck dan Optimasi}

\textit{Bottleneck} utama sistem bukan pada komputasi CPU melainkan pada \textit{network latency}: (a) menunggu halaman web BI memuat, (b) menunggu antrean API MinerU, dan (c) jeda anti-bot yang disengaja. Hal ini sesuai dengan karakteristik umum pipeline scraping web, di mana waktu tunggu jaringan mendominasi keseluruhan waktu eksekusi.

Penggunaan parameter \texttt{max\_active\_tis\_per\_dag=2} yang membatasi dua Chrome instance berjalan bersamaan terbukti efektif mencegah kelebihan beban memori (\textit{memory exhaustion}). Untuk sistem dengan RAM lebih besar atau jumlah worker yang lebih banyak, nilai ini dapat ditingkatkan secara proporsional.

% ============================================================
\section{KESIMPULAN}
% ============================================================

Penelitian ini berhasil merancang dan mengimplementasikan sebuah \textit{pipeline data engineering} terdistribusi yang secara otomatis mengekstrak indikator kesenjangan ekonomi regional Indonesia dari Laporan Perekonomian Provinsi Bank Indonesia. Sistem mengintegrasikan empat komponen utama: Selenium untuk \textit{web scraping} adaptif, MinerU VLM untuk ekstraksi PDF berbasis AI, Apache Airflow untuk orkestrasi dan penjadwalan, serta MinIO untuk penyimpanan berbasis \textit{Data Lake} dengan arsitektur Medallion.

Eksperimen pada lingkungan dua node menunjukkan bahwa distribusi beban komputasi berhasil dilakukan secara otomatis melalui Celery Executor dan Redis, dengan node Master memproses 22 provinsi dan node Worker menyelesaikan 6 provinsi secara paralel. Penerapan \textit{Dynamic Task Mapping} pada Airflow memungkinkan satu definisi task menjadi 34 instance independen, masing-masing bertanggung jawab atas satu provinsi. Mekanisme \textit{auto-retry} berhasil menangani kegagalan sementara akibat ketidakstabilan jaringan dan \textit{timeout} API. \textit{Silver Layer} yang dihasilkan berisi tujuh dari delapan file CSV indikator ekonomi yang siap untuk analisis lebih lanjut menggunakan alat visualisasi seperti Metabase atau Tableau.

% ============================================================
% REFERENSI (IEEE Style)
% ============================================================
\begin{thebibliography}{99}

\bibitem{bps2023}
Badan Pusat Statistik, ``Distribusi Pendapatan Penduduk Indonesia, Koefisien Gini 2023,'' BPS, Jakarta, 2023. [Online]. Available: \url{https://www.bps.go.id}. [Accessed: 07-Jun-2026].

\bibitem{py_pdf_extraction}
S. Bhatt and P. Bhatt, ``PDF Data Extraction Using Machine Learning: A Survey,'' \textit{International Journal of Advanced Research in Computer Science}, vol. 11, no. 3, pp. 45--51, 2020.

\bibitem{mineru2024}
Y. Wang, S. Lv, L. Huang, et al., ``MinerU: An Open-Source Solution for Precise Document Content Extraction,'' \textit{arXiv preprint arXiv:2409.18839}, 2024.

\bibitem{harenslak2021}
B. Harenslak and J. de Ruiter, \textit{Data Pipelines with Apache Airflow}. Shelter Island, NY: Manning Publications, 2021.

\bibitem{databricks_medallion}
Databricks Inc., ``What is a Medallion Architecture?,'' Databricks Documentation, 2023. [Online]. Available: \url{https://www.databricks.com/glossary/medallion-architecture}. [Accessed: 07-Jun-2026].

\bibitem{airflow_docs}
Apache Software Foundation, ``Apache Airflow Documentation: Celery Executor,'' 2024. [Online]. Available: \url{https://airflow.apache.org/docs/apache-airflow/stable/core-concepts/executor/celery.html}. [Accessed: 07-Jun-2026].

\bibitem{zaharia2016}
M. Zaharia, R. S. Xin, P. Wendell, et al., ``Apache Spark: A Unified Engine for Big Data Processing,'' \textit{Communications of the ACM}, vol. 59, no. 11, pp. 56--65, Nov. 2016.

\bibitem{inmon2021}
W. H. Inmon, ``Data Lakehouse: A New Paradigm for Data Management,'' in \textit{Proceedings of the International Conference on Data Engineering}, 2021, pp. 1--15.

\end{thebibliography}

\end{document}